\documentclass{article}

\usepackage[T1]{fontenc}
\usepackage{lmodern}
\usepackage{amsmath}
\usepackage{amssymb}
\usepackage{amsthm}
\usepackage{mathtools}
\usepackage[hidelinks]{hyperref}

\newtheorem{proposition}{Proposition}

\DeclareMathOperator{\wt}{wt}
\DeclareMathOperator{\supp}{supp}
\DeclareMathOperator{\row}{row}

\title{A column criterion for CSS distance at least three}
\author{}
\date{Standalone analytical draft, 11 November 2026}

\begin{document}

\maketitle

\textbf{Status: PROPOSITION throughout; independent mathematical review pending.}
The Lean module and guarded axiom audit passed in the coordinating capped detached run, as
recorded at the end. The arguments below are complete proposed proofs for hand checking,
not final mathematically reviewed proofs. No concrete 722-code distance is asserted.

\section{Main column criterion}\label{sec:criterion}

\begin{proposition}\label{prop:criterion}
Let \(F = \mathbb{F}_2 = \{0, 1\}\), with addition and multiplication modulo two. Let \(Q\)
be a finite set of physical coordinates and \(R\) any set of check indices. Let
\(H \colon R \times Q \to F\) be a matrix and \(B\) a linear subspace of \(F^Q\). For
\(i \in Q\), define the unit vector \(e_i\) by \(e_i(i) = 1\) and \(e_i(q) = 0\) for
\(q \ne i\). Define the column \(c_i \in F^R\) and the syndrome \(Hx\) by
\begin{equation}
  c_i(r) = H(r, i), \quad (Hx)(r) = \sum_{q \in Q} H(r, q) x(q).
\end{equation}
For \(x \in F^Q\), its support and Hamming weight are
\begin{equation}
  \supp(x) = \{q \in Q : x(q) \ne 0\}, \quad \wt(x) = \lvert\supp(x)\rvert.
\end{equation}
Assume the following two hypotheses:
\begin{enumerate}
  \item \textbf{No zero columns:} \(c_i \ne 0\) for every \(i \in Q\).
  \item \textbf{Trivial column collisions:} for every distinct \(i, j \in Q\), if
    \(c_i = c_j\), then \(e_i + e_j \in B\).
\end{enumerate}
Then
\begin{equation}
  Hx = 0 \text{ and } \wt(x) \le 2 \quad \Longrightarrow \quad x \in B.
\end{equation}
Consequently every \(x\) with \(Hx = 0\) and \(x \notin B\) satisfies \(\wt(x) \ge 3\). No
inclusion \(B \subseteq \ker H\) is needed for this implication. Such an inclusion will
hold in the CSS applications below.
\end{proposition}

\begin{proof}[Proposed proof (PROPOSITION; review pending)]
  A vector over \(F\) is exactly the indicator of its support: at a supported coordinate
  its nonzero value must be 1, and at every other coordinate it is 0. A finite set of
  cardinality at most two is empty, a singleton, or a pair of distinct elements. This
  gives the following exhaustive classification, with no enumeration of \(Q\):
  \begin{itemize}
    \item If \(\supp(x)\) is empty, every coordinate is zero, so \(x = 0\).
    \item If \(\supp(x) = \{i\}\), the indicator description gives \(x = e_i\).
    \item If \(\supp(x) = \{i, j\}\) with \(i \ne j\), the indicator description gives
      \(x = e_i + e_j\). At \(i\) and \(j\) exactly one summand is 1; elsewhere both are 0.
  \end{itemize}
  In the first case \(x \in B\) because a linear subspace contains zero. In the second
  case, the definition of matrix multiplication gives \(He_i = c_i\); thus \(Hx = 0\)
  contradicts the first hypothesis. In the third case, linearity gives
  \(0 = Hx = c_i + c_j\). Adding \(c_j\) to both sides and using \(c_j + c_j = 0\) gives
  \(c_i = c_j\). The second hypothesis now puts \(x = e_i + e_j\) in \(B\). This proves the
  first conclusion. If \(Hx = 0\) and \(x \notin B\), a weight at most two would contradict
  that conclusion. Weight is a natural number, so it is at least three. This completes the
  proposed argument.
\end{proof}

The statement allows repeated columns and dependent check rows. It requires neither a
positive-dimensional logical space nor an existing nontrivial cycle. If there are no such
cycles, the universal lower bound is vacuous.

\section{Definitions for the CSS specialization}\label{sec:css-definitions}

Take finite check sets \(R_X, R_Z\) and binary matrices \(H_X \colon R_X \times Q \to F\)
and \(H_Z \colon R_Z \times Q \to F\) satisfying
\begin{equation}
  H_X H_Z^T = 0.
\end{equation}
Here transpose interchanges row and column indices. This equation says that every X-check
row has zero dot product with every Z-check row, where
\(u \cdot v = \sum_{q \in Q} u(q) v(q)\) in \(F\). Define the stabilizer subspaces and
cycle spaces by
\begin{gather}
  U = \row(H_X), \quad V = \row(H_Z), \\
  C_X = \ker H_Z, \quad C_Z = \ker H_X.
\end{gather}
The row space is the set of all binary linear combinations of rows; no row independence is
assumed. An X-sector nontrivial logical is an \(x \in C_X\) with \(x \notin U\). A Z-sector
nontrivial logical is a \(z \in C_Z\) with \(z \notin V\). Thus the X sector uses the
Z-check syndrome, and conversely.

\section{CSS sector corollaries}\label{sec:css}

\begin{proposition}\label{prop:css}
For the X-sector bound, assume:
\begin{enumerate}
  \item Every column of \(H_Z\) is nonzero.
  \item Whenever distinct columns \(i, j\) of \(H_Z\) are equal, \(e_i + e_j \in U\).
\end{enumerate}
Then every nontrivial X-sector logical has Hamming weight at least three. For the Z-sector
bound, assume:
\begin{enumerate}
  \item Every column of \(H_X\) is nonzero.
  \item Whenever distinct columns \(i, j\) of \(H_X\) are equal, \(e_i + e_j \in V\).
\end{enumerate}
Then every nontrivial Z-sector logical has Hamming weight at least three. Either
conclusion requires only its own two column hypotheses. If both sets of hypotheses hold,
both conclusions hold simultaneously.
\end{proposition}

\begin{proof}[Proposed proof (PROPOSITION; review pending)]
  First, commutation makes the stabilizer spaces subspaces of the corresponding cycle
  spaces. Indeed, if \(u = \sum_{r \in R_X} a_r H_X(r, \cdot)\), then for every
  \(s \in R_Z\),
  \begin{equation}
    (H_Z u)(s) = \sum_{r \in R_X} a_r
      \Bigl(\sum_{q \in Q} H_Z(s, q) H_X(r, q)\Bigr) = 0.
  \end{equation}
  The inner sum is an entry of \(H_X H_Z^T\), using commutativity of multiplication in
  \(F\). Hence \(U \subseteq C_X\). Likewise, for \(v = \sum_{s \in R_Z} b_s H_Z(s, \cdot)\),
  each coordinate of \(H_X v\) is a linear combination of entries of \(H_X H_Z^T\), so
  \(V \subseteq C_Z\).

  Now apply Proposition~\ref{prop:criterion} with \((H, B) = (H_Z, U)\). Its hypotheses are
  precisely the two X-sector hypotheses. For \(x \in C_X\) outside \(U\) it gives
  \(\wt(x) \ge 3\). Apply it with \((H, B) = (H_X, V)\) for the Z sector. This proves both
  claimed bounds.
\end{proof}

\textbf{Distance convention.} If the set \(C_X \setminus U\) is nonempty, define \(d_X\) as
the minimum Hamming weight of its elements, and similarly define \(d_Z\) when
\(C_Z \setminus V\) is nonempty. The preceding universal statements then give
\(d_X \ge 3\) and \(d_Z \ge 3\) whenever the respective minimum exists. No numerical value
is assigned here to an empty minimum. In particular the corollary does not assert that a
code encodes a logical qubit or that its distance equals three.

\section{Definitions for mixed Pauli operators}\label{sec:pauli-definitions}

On one qubit, use the matrices
\begin{equation}
  I = \begin{pmatrix} 1 & 0 \\ 0 & 1 \end{pmatrix}, \quad
  X = \begin{pmatrix} 0 & 1 \\ 1 & 0 \end{pmatrix}, \quad
  Z = \begin{pmatrix} 1 & 0 \\ 0 & -1 \end{pmatrix}.
\end{equation}
A Pauli operator on \(Q\) is, up to a scalar phase in
\(\{1, -1, \mathrm{i}, -\mathrm{i}\}\), specified by a pair
\((x, z) \in F^Q \times F^Q\): at coordinate \(q\) its operator is
\(X^{x(q)} Z^{z(q)}\). Tensor products over all coordinates give the full operator. The
pair \((0, 0)\) is identity at that coordinate; \((1, 0)\), \((0, 1)\), and \((1, 1)\) are
its three nonidentity possibilities up to phase. Define
\begin{equation}
  \wt_P(x, z) = \lvert\supp(x) \cup \supp(z)\rvert.
\end{equation}
Global phase does not change the pair or this weight. In particular
\begin{equation}
  \wt_P(x, z) \ge \wt(x), \quad \wt_P(x, z) \ge \wt(z).
\end{equation}
These inequalities follow from the two inclusions into the union and monotonicity of
finite-set cardinality.

For the CSS matrices of Proposition~\ref{prop:css}, let \(S\) be the group generated by the
pure-X check operators \((H_X(r, \cdot), 0)\) and pure-Z check operators
\((0, H_Z(s, \cdot))\), with their standard phase-one representatives. A nontrivial Pauli
logical means an operator commuting with every element of \(S\) whose operator pair is not
the pair of any element of \(S\). Thus a scalar multiple of a stabilizer is excluded,
including a phase-only operator. This is the convention needed for code distance.

\section{Mixed-Pauli corollary}\label{sec:pauli}

\begin{proposition}\label{prop:pauli}
Under both sets of column hypotheses in Proposition~\ref{prop:css}, every nontrivial Pauli
logical for \(S\) has Pauli weight at least three.
\end{proposition}

\begin{proof}[Proposed proof (PROPOSITION; review pending)]
  Direct multiplication gives \(X^2 = Z^2 = I\) and \(ZX = -XZ\). Operators on different
  coordinates commute. Consequently the commutation sign between pairs \((x, z)\) and
  \((u, v)\) is
  \begin{equation}
    (-1)^{x \cdot v + z \cdot u}.
  \end{equation}
  Indeed, moving every X factor through the other operator's Z factors introduces exactly
  one minus sign for each common coordinate with both bits 1; all exponents are summed
  modulo two. Products of Pauli pairs add their two binary components, regardless of the
  resulting scalar phase.

  The CSS equation therefore makes all the check generators commute. A pair \((x, z)\)
  commutes with every X generator exactly when \(H_X z = 0\), and with every Z generator
  exactly when \(H_Z x = 0\). Commuting with generators is equivalent to commuting with the
  generated group: commutation is preserved under products and inverses. Thus a Pauli
  logical has \(x \in C_X\), \(z \in C_Z\).

  The set of operator pairs of \(S\) is exactly \(U \times V\). One inclusion follows by
  adding the pairs of the generators in a product. For the reverse inclusion, express
  \(u \in U\) as a sum of X rows and \(v \in V\) as a sum of Z rows, and multiply the
  selected generators; the resulting pair is \((u, v)\). Global phases do not affect this
  statement. In fact these commuting generators define a valid stabilizer group: reorder
  each product with all X generators first, giving the phase-one representative
  \(X^u Z^v\); if its pair is zero, this representative is \(I\), so \(-I\) is not
  generated.

  Nontriviality therefore means \(x \notin U\) or \(z \notin V\). In the first case
  Proposition~\ref{prop:css} gives \(\wt(x) \ge 3\), so \(\wt_P(x, z) \ge 3\). In the second
  case it gives \(\wt(z) \ge 3\) and the same conclusion. This proves the corollary without
  assuming that the logical is purely X or purely Z.
\end{proof}

\section{Explicit hypotheses for a possible 722 adapter}\label{sec:adapter}

This section specifies an application interface only. It does not establish that any
submitted or reconstructed 722-qubit matrices meet that interface. In particular the
symbolic column classification in another draft is not silently taken as a proved premise
here.

Take \(Q = \{0, \dots, 721\}\) and name its qubits
\begin{equation}
  A(i, j) = 19i + j, \quad B(i, j) = 361 + 19i + j \quad (0 \le i, j \le 18).
\end{equation}
Suppose a CSS code has matrices \(H_X, H_Z\), each with 357 rows, and has the following
verified properties:
\begin{enumerate}
  \item \textbf{Input identification and commutation.} The matrices are identified with
    the intended literal input on these physical coordinates, and \(H_X H_Z^T = 0\). X rows
    are labeled \(X_{c,d}\) with \(-2 \le c \le 18\), \(0 \le d \le 16\); Z rows are
    labeled \(Z_{c,d}\) with \(2 \le c \le 18\), \(0 \le d \le 20\). The labels must refer
    to actual rows of the recognized matrices, not to a different board or boundary
    convention.
  \item \textbf{Nonzero Z-check columns.} Every column of \(H_Z\) is nonzero.
  \item \textbf{Complete Z-check collision classification.} If distinct columns of \(H_Z\)
    are equal, their qubits, in either order, are \(A(18, j), B(18, j - 2)\) for some
    integer \(2 \le j \le 18\).
  \item \textbf{X stabilizer witnesses.} For every such \(j\), the actual row
    \(X_{18,j-2}\) is \(e_{A(18,j)} + e_{B(18,j-2)}\).
  \item \textbf{Nonzero X-check columns.} Every column of \(H_X\) is nonzero.
  \item \textbf{Complete X-check collision classification.} If distinct columns of \(H_X\)
    are equal, their qubits, in either order, are \(A(i, 18), B(i + 2, 18)\) for some
    integer \(0 \le i \le 16\).
  \item \textbf{Z stabilizer witnesses.} For every such \(i\), the actual row
    \(Z_{i+2,20}\) is \(e_{A(i,18)} + e_{B(i+2,18)}\).
\end{enumerate}

\begin{proposition}[conditional adapter; review pending]\label{prop:adapter}
  If all these hypotheses hold, the universal X-sector, Z-sector, and mixed-Pauli lower
  bounds of three hold for that recognized code.
\end{proposition}

\begin{proof}[Proposed proof (PROPOSITION; review pending)]
  For any Z-check column collision, the classification and X witness put the two-site
  vector in \(\row(H_X)\): an individual row is a linear combination of rows with its own
  coefficient 1 and all others 0. Exchanging the two sites does not change their sum. The
  same argument with the X-check classification and Z witness puts each X-check collision
  in \(\row(H_Z)\). Together with the nonzero-column hypotheses these are exactly both
  instances of Proposition~\ref{prop:css}. Apply Propositions~\ref{prop:css}
  and~\ref{prop:pauli}. This implication does not discharge any of its premises.
\end{proof}

\section{Formalization correspondence and review boundary}\label{sec:formalization}

\begin{sloppypar}
The companion draft is
\nolinkurl{lean/QldpcVerification/Planar/SmallSupport.lean}.
Its \texttt{support} and \nolinkurl{card_support} use physical support and Mathlib's
\texttt{hammingNorm}; \nolinkurl{classify_weight_le_two} uses
\nolinkurl{Finset.card_eq_one} and \nolinkurl{Finset.card_eq_two}. The generic conclusions
are \nolinkurl{small_cycle_mem} and \nolinkurl{weight_ge_three}. The two hypotheses are
packaged as \texttt{ColumnCriterion}. The project CSS corollaries are \nolinkurl{x_lower}
and \nolinkurl{z_lower}, using \nolinkurl{Input.CSSCode.IsXLogical}, \texttt{IsZLogical},
\texttt{xStab}, and \texttt{zStab}. The theorem \nolinkurl{pauli_lower} transfers through
the project's homological identification and QECLean's
\nolinkurl{chainWeight_lower_bound_transfers}.
\end{sloppypar}

\begin{sloppypar}
The separate file
\nolinkurl{lean/QldpcVerification/PlanarSmallSupportAudit.lean}
contains guarded axiom checks. These filenames identify draft interfaces, not substitutes
for mathematical review. The coordinating run reported that both the module and guarded
audit passed in the capped detached allocation \texttt{small-v2} (3.841 s and 2.807 s;
reported process peak 3,774,580 KiB). This is a reported formal check, not an independent
mathematical review. The actual 722 adapter remains separate and unproved in this
document. No support-search certificate, numerical distance estimate, rank computation, or
assertion of distance twenty is a premise of this note.
\end{sloppypar}

\end{document}
